\section*{Chapitre \ref{ch:g9:metals-acids-corrosion} --- Les métaux~: réactions avec les acides et corrosion}

\begin{solution}{exo:g9:metals-acids-corrosion:1}
Du dihydrogène~: une bûchette enflammée à l'ouverture du tube produit un
«~pop~» aigu.
\end{solution}

\begin{solution}{exo:g9:metals-acids-corrosion:2}
Le fer, le zinc et l'aluminium réagissent~; le cuivre et l'or, non.
\end{solution}

\begin{solution}{exo:g9:metals-acids-corrosion:3}
Le dioxygène et l'eau, tous les deux.
\end{solution}

\begin{solution}{exo:g9:metals-acids-corrosion:4}
Un \omterm{def:g9:ions:ion}{ion} présent pendant une réaction mais qui n'y participe pas. Avec
l'acide chlorhydrique~: l'\omterm{def:g9:ions:ion}{ion} chlorure \ce{Cl-}.
\end{solution}

\begin{solution}{exo:g9:metals-acids-corrosion:5}
\ce{Mg(s) + 2H+(aq) -> Mg^{2+}(aq) + H2(g)}~; charge $+2$ de chaque côté.
\end{solution}

\begin{solution}{exo:g9:metals-acids-corrosion:6}
Ajouter de l'hydroxyde de sodium~: un \omterm{def:g9:ions:precipitate}{précipité} vert de \ce{Fe(OH)2}
révèle les \omterm{def:g9:ions:ion}{ions} fer(II).
\end{solution}

\begin{solution}{exo:g9:metals-acids-corrosion:7}
L'ébullition chasse l'\omterm{def:g4:air-a-mixture-of-gases:air}{air} dissous dans l'eau, et l'huile empêche l'\omterm{def:g4:air-a-mixture-of-gases:air}{air}
de s'y dissoudre de nouveau~: le clou a de l'eau mais pas de dioxygène.
\end{solution}

\begin{solution}{exo:g9:metals-acids-corrosion:8}
Le sel dissous dans l'eau projetée sur la voiture accélère la \omterm{def:g9:metals-acids-corrosion:corrosion}{rouille}.
\end{solution}

\begin{solution}{exo:g9:metals-acids-corrosion:9}
La couche d'oxyde formée adhère au métal et l'isole de l'\omterm{def:g4:air-a-mixture-of-gases:air}{air}~: l'attaque
s'arrête à la surface.
\end{solution}

\begin{solution}{exo:g9:metals-acids-corrosion:10}
Le peindre, l'huiler ou le graisser, le galvaniser (ou le fabriquer en
acier inoxydable).
\end{solution}

\begin{solution}{exo:g9:metals-acids-corrosion:11}
\ce{2Al(s) + 6H+(aq) -> 2Al^{3+}(aq) + 3H2(g)}~: 2 Al et 6 H de chaque
côté, charge $+6$ de chaque côté. Pour 200 \omterm{def:g7:atoms-and-molecules:atom}{atomes} d'Al~: $200 \times \frac{3}{2} = 300$
\omterm{def:g7:atoms-and-molecules:molecule}{molécules} de dihydrogène.
\end{solution}

\begin{solution}{exo:g9:metals-acids-corrosion:12}
Le zinc autour de la rayure se corrode à la place de l'acier~: tant
qu'il reste du zinc à proximité, l'acier est épargné.
\end{solution}

\begin{solution}{pb:g9:metals-acids-corrosion:1}
\textbf{1.} \ce{Zn(s) + 2H+(aq) -> Zn^{2+}(aq) + H2(g)}.

\textbf{2.} Le dihydrogène~: un «~pop~» aigu avec une bûchette
enflammée.

\textbf{3.} L'hydroxyde de sodium donne un \omterm{def:g9:ions:precipitate}{précipité} blanc de
\ce{Zn(OH)2}.

\textbf{4.} Les \omterm{def:g9:ions:ion}{ions} chlorure \ce{Cl-}.

\textbf{5.} $125.6 - 113.5 = \qty{12.1}{g}$.

\textbf{6.} Du dihydrogène se dégage tant que le zinc réagit avec
l'acide~; quand il ne reste plus de zinc, l'effervescence s'arrête.

\textbf{7.} Le fer de l'acier réagirait lui aussi, en donnant du
dihydrogène et des \omterm{def:g9:ions:ion}{ions} \ce{Fe^{2+}}~; l'hydroxyde de sodium donnerait
alors un \omterm{def:g9:ions:precipitate}{précipité} vert.

\textbf{8.} Le zinc est attaqué à la place de l'acier et tient l'eau et
l'\omterm{def:g4:air-a-mixture-of-gases:air}{air} à distance de celui-ci.

\textbf{9.} $2 \times 10.0 \times 10.0 = \qty{200}{cm^2}$.

\textbf{10.} $12.1 \div 7.13 \approx \qty{1.70}{cm^3}$.

\textbf{11.} $1.70 \div 200 = \qty{0.0085}{cm}$.

\textbf{12.} $0.0085 \times 10\,000 = \qty{85}{\micro m}$~: la couche de
zinc avait une épaisseur d'environ \qty{85}{\micro m}.
\end{solution}
